A popular prescription for automated scientific ideation is to \emph{ground} the
process---in deduction, induction, small experiments, literature comparison, or
proposal writing---rather than prompting a language model abstractly. Yet every
prior method grounds in only one way and compares against a (typically weak)
ungrounded baseline, so it remains unknown which grounding \emph{mechanism}
actually helps when all else is held fixed. We present the first single-protocol,
head-to-head ablation of six ideation regimes (one ungrounded control plus five
grounding mechanisms). Holding the model (\gptmini), task suite, generation
budget, and evaluation procedure fixed, each regime produces a bank of $15$ task
hypotheses, and each bank is scored by an objective downstream-utility metric: how
accurately an LLM applying the bank predicts held-out labels, both in-distribution
(\IND) and out-of-distribution (\OOD), across three classification tasks and three
seeds ($54$ banks, $\approx 11.9$k API calls). We find that grounding does
\emph{not} uniformly help: pooled across regimes, grounded banks are statistically
indistinguishable from---and slightly below---the ungrounded control. But the
mechanism matters decisively. Grounding in \emph{empirical feedback}
(test-and-revise) is the only regime that significantly improves \IND{} accuracy
($+4.4$ points, Cohen's $d=1.09$, $p=0.02$), at roughly $80\times$ the token cost;
proposal writing significantly \emph{hurts} ($-5.7$ points, $p=0.008$) and
collapses idea diversity; and data grounding can \emph{overfit}, trading \OOD{}
robustness for \IND{} fit. Novelty and diversity do not track usefulness within a
task. For today's strong LLMs, ``ground your ideation'' is too coarse: only
empirical grounding reliably adds value, and ungrounded prompting is a
surprisingly hard baseline to beat.
